\documentclass[11pt]{article}
\usepackage{amsmath,amssymb,amsthm,mathtools}
\usepackage[margin=1in]{geometry}
\usepackage{booktabs,array,enumitem,xcolor}
\usepackage[hidelinks]{hyperref}

\newtheorem{theorem}{Theorem}[section]
\newtheorem{proposition}[theorem]{Proposition}
\newtheorem{lemma}[theorem]{Lemma}
\newtheorem{corollary}[theorem]{Corollary}
\newtheorem{definition}[theorem]{Definition}
\newtheorem{conjecture}[theorem]{Conjecture}
\theoremstyle{remark}
\newtheorem{remark}[theorem]{Remark}

\newcommand{\R}{\mathbb R}
\newcommand{\Q}{\mathbb Q}
\newcommand{\Z}{\mathbb Z}
\newcommand{\KG}{K_G}
\newcommand{\kc}{\kappa_c}
\newcommand{\ks}{\kappa_s}
\newcommand{\kint}{\kappa_{\mathrm{int}}}
\newcommand{\Cham}{\mathcal C}
\newcommand{\Wp}{W_+}
\newcommand{\rank}{\operatorname{rank}}
\newcommand{\Gal}{\operatorname{Gal}}
\newcommand{\status}[1]{\textup{\textsc{#1}}}
\newcommand{\established}{\status{Established}}
\newcommand{\realization}{\status{Realization}}
\newcommand{\bridge}{\status{Bridge}}
\newcommand{\conditional}{\status{Conditional}}
\newcommand{\proofstop}[1]{\par\medskip\noindent{\color{red!75!black}\textbf{STOP--PROOF.}} \textit{#1}\par\medskip}

\title{\textbf{Active Recharting of Constraint Hypersurfaces}\\
\large Orbit Quotients, Channel Accounts, and Physical Selection\\[4pt]
\large A local-to-global synthesis with a Galois--spectral realization}
\author{William Lloyd}
\date{Draft v0.3 --- 13 July 2026}

\begin{document}
\maketitle

\begin{abstract}
We assemble a local and a global branch into one theorem spine. Locally, a
labelled three-coordinate constraint hypersurface is represented by its active
output-chart jet orbit: a functional is invariant under active recharting
exactly when it factors through the active orbit quotient. Curvature channel
data forms a framed, defining-function-relative account whose canonical sum
is an elementary shape-operator invariant. In the Kerr--Newman realization,
positivity and interior regularity select a unique member of a stated
inverse-channel metric family, and local metric germs classify its open
boundary strata.

Globally, the four-charge seed-mass relation gives a finite signed-sheet
cover. The axial k-ellipse pencil identifies the strict mass chamber with a
definiteness chamber. An involutive parity block decomposition gives an exact
singular-value factorization of the mass-square polynomial; every root is
positive, every positive-imbalance sheet crosses once and transversely, and
the physical root is the unique least root and first spectrahedral exit.

The unification claimed here is a theorem dependency and comparison, not an
identification of group actions. The active chart group, radical sign group,
axial deck involution, and global monodromy act on different carriers. Every
point at which an identification, a generic-to-specialized arithmetic
promotion, or a certificate-only calculation would be required is marked
STOP--PROOF. The result is both a synthesis and an explicit proof programme.
\end{abstract}

\tableofcontents

\section{Introduction}

\subsection{Scope and terminology}

The abstract local datum in this paper is a constraint hypersurface equipped
with a \emph{labelled coordinate frame}. Its formal structure is the ability
to use each coordinate as an output chart and compare the resulting finite
jets under active recharting. In a directed application, the three coordinates
may be interpreted as Directive, Substrate, and Product; we call that a
\emph{DBP realization}. Those semantics are not used in the abstract
definitions or theorem names. This distinction is essential because the global
Galois--spectral branch has no natural Directive--Substrate--Product
interpretation.

\subsection{The spine}

The common logical order is
\[
\boxed{\begin{gathered}
\text{carrier}\longrightarrow\text{orbit quotient}
\longrightarrow\text{account fibre}\\
\longrightarrow\text{admissible selection}
\longrightarrow\text{typed strata}
\longrightarrow\text{local normal form}
\end{gathered}}
\]
The local and global branches instantiate this order differently:
\[
\begin{array}{c@{\qquad}c}
\textbf{framed constraint geometry}&\textbf{global cover geometry}\\[1mm]
\text{active output-chart atlas}&\text{signed radical carrier}\\
\downarrow&\downarrow\\
\text{chart-orbit quotient}&\text{norm eliminant and }U=m^2\\
\downarrow&\downarrow\\
\text{curvature-channel account}&\text{physical/shadow root fibre}\\
\downarrow&\downarrow\\
\text{admissible metric}&\text{first definite-pencil exit}\\
\downarrow&\downarrow\\
\text{chart/isotropy strata}&\text{collision/wall strata}.
\end{array}
\]
The central assertion is not that the columns are equal. Both exhibit a
rigorous separation between invariant reduction and physical selection.

\subsection{Status discipline}

\begin{center}
\begin{tabular}{>{\bfseries}p{0.18\textwidth}p{0.70\textwidth}}
\toprule
Established & A self-contained proof is included or reduces to a precise
classical theorem whose hypotheses are checked.\\
Realization & Proved for a named family, not promoted to a universal statement
for framed constraint hypersurfaces.\\
Bridge & An exact comparison or common architecture; no equivalence is
asserted.\\
Conditional & Depends on a banked computation, certificate, or companion
proof that must be re-verified before publication.\\
\bottomrule
\end{tabular}
\end{center}

\proofstop{The word unified means an ordered theorem spine. Promoting
the spine to an autonomous category requires native morphisms, equivalence,
closure, and an independent third realization; none is supplied here.}

\section{Framed constraint hypersurfaces and non-identification rules}

\subsection{The local chart datum}

Let $F(x_1,x_2,x_3)=0$ be a smooth local constraint relation. On
\[
\mathcal U_1=\{F=0:F_{x_1}F_{x_2}F_{x_3}\neq0\},
\]
the relation may be solved with each coordinate as output. Fix $r\geq1$,
write $X_r$ for the regular finite-jet space, and $J_r(x)$ for one graph jet.

\begin{definition}[active chart-jet carrier]
The group $G_{\mathrm{act}}=S_3$ acts by re-solving the same relation in
each output coordinate and transporting the jet by the finite inverse-function
theorem. The carrier is
\[
\mathcal O_r(x)=G_{\mathrm{act}}\cdot J_r(x).
\]
Retaining coordinate values gives
\[
\mathcal O_r^+(x)
=G_{\mathrm{act}}\cdot(x_1,x_2,x_3,J_r(x)).
\]
\end{definition}

\subsection{Global cover data}

For real parameters $(M,N_1,\ldots,N_k)$ define
\[
w_i(U)=\sqrt{U+N_i^2},\qquad
s_\varepsilon(U)=\sum_i\varepsilon_iw_i(U),
\quad\varepsilon\in\{\pm1\}^k.
\]
The signed norm is the polynomial
\[
N_k(U)=\prod_{\varepsilon\in\{\pm1\}^k}
\left(4M-s_\varepsilon(U)\right),
\]
and its oriented axial lift is $q_k(m)=N_k(m^2)$. We distinguish
\[
H_{\mathrm{sign}}=(C_2)^k,\qquad
G_{\mathrm{deck}}=\langle\tau\rangle,\quad\tau(m)=-m,
\]
from $G_{\mathrm{act}}$ and any global monodromy group; here $C_2$ is the
cyclic group of order two.

\begin{proposition}[non-identification guardrail]
The constructions do not provide canonical identifications
\[
S_3\not\cong(C_2)^k,\qquad
X_r/S_3\not\cong\operatorname{Spec}\Q[U]/(N_k),\qquad
\Sigma_r\not\equiv\operatorname{Norm}.
\]
Local chart recharting and global analytic-continuation monodromy also act on
different sets.
\end{proposition}

\begin{proof}
The chart action permutes output charts of one local relation. Radical sign
changes move signed sheets, and the deck involution changes an axial cover
coordinate. No map intertwining these actions has been defined.
\end{proof}

\proofstop{Any stronger claim requires an explicit carrier map, proof of
equivariance, and comparison of stabilisers, branch loci, and quotients. A
structural resemblance is not such a proof.}

\section{The Active Recharting Orbit Theorem}

\begin{definition}[active-recharting invariant]
A function $\phi:X_r\to Y$ is an order-$r$ active-recharting invariant if
\[
\phi(\mathcal T_\sigma\xi)=\phi(\xi)
\]
for every $\sigma\in S_3$ wherever both sides are defined.
\end{definition}

\begin{theorem}[Active Recharting Orbit Theorem]\label{thm:orbit}
Let $\pi_r:X_r\to X_r/S_3$ be the set-orbit map. Then
$\phi:X_r\to Y$ is an order-$r$ active-recharting invariant if and only if
there is a unique $\bar\phi:X_r/S_3\to Y$ satisfying
\[
\phi=\bar\phi\circ\pi_r.
\]
The same holds for the value-plus-jet carrier $X_r^+$.
\end{theorem}

\begin{proof}
The active transition maps satisfy the $S_3$ relations on their common
regular domain. A function is invariant exactly when it is constant on
orbits, which is equivalent to unique factorization through the orbit set.
\end{proof}

\begin{proposition}[rational finite-jet action]\label{prop:rationalaction}
At orders one and two, active output-chart recharting is a rational $S_3$
action whose
denominators are powers of the relevant first derivatives. At each fixed
finite order the same conclusion follows recursively from the finite
inverse-function theorem.
\end{proposition}

\proofstop{For an all-$r$ headline, include the germ-level composition
argument and a Fa\'{a}-di-Bruno denominator induction. Otherwise restrict
the published theorem to $r\leq2$, where all formulas are explicit.}

\begin{remark}[mixed invariants]
A derivative-only tuple is not automatically complete for invariants mixing
values and jets. Such invariants belong on $X_r^+$. The orbit theorem avoids
claiming that a preferred finite kernel is complete without proving orbit
separation.
\end{remark}

\section{Channel accounts and invariant reduction}

Let $F:\R^n\to\R$ define a regular hypersurface point and put
\[
g=\nabla F,\qquad H=\nabla^2F,\qquad q=g^{\mathsf T}g\neq0.
\]
In the distinguished coordinate frame, split
\[
H=H_c+H_s,\qquad H_s=\operatorname{diag}(H),\qquad H_c=H-H_s.
\]
For $1\leq r\leq n-1$ define
\[
\widehat{\mathcal C}_r(t,v)
=(-1)^{r+1}\sum_{|I|=r+1}
\det\begin{pmatrix}
0&g_I^{\mathsf T}\\
g_I&(tH_c+vH_s)_I
\end{pmatrix}
=\sum_{p+s=r}t^pv^s\widehat\kappa_{r;p,s}.
\]
The letter $v$ avoids collision with the mass-square coordinate $U=m^2$.
Set
\[
\kappa_{r;p,s}
=\frac{\widehat\kappa_{r;p,s}}{q^{(r+2)/2}}.
\]

\begin{theorem}[bordered-minor channel reduction]\label{thm:reduction}
Fix the oriented shape operator
\[
P=I-\frac{gg^{\mathsf T}}q,\qquad
S=-\frac{PHP}{\sqrt q}.
\]
Then
\[
\sigma_r(S)
=\frac{\widehat{\mathcal C}_r(1,1)}{q^{(r+2)/2}}
=\sum_{p+s=r}\kappa_{r;p,s}.
\]
If $V_r=\R^{r+1}$ is the ambient coefficient space and
\[
\Sigma_r:V_r\to\R,\qquad
(c_0,\ldots,c_r)\longmapsto\sum_jc_j,
\]
then
\[
0\longrightarrow\ker\Sigma_r
\longrightarrow V_r
\mathop{\longrightarrow}^{\Sigma_r}\R
\longrightarrow0
\]
is exact. A geometrically realizable channel vector lies in the affine
hyperplane over the reduced elementary extrinsic curvature $\sigma_r$.
\end{theorem}

\begin{proof}[Proof route]
Rotate $g$ to $\sqrt q\,e_n$. Only bordered minors containing the normal
coordinate survive; each is $-q$ times the corresponding tangent principal
minor. The sign and normalization give $\sigma_r(S)$. Replacing $H$ by
$tH_c+vH_s$ gives a homogeneous polynomial, and evaluation at $(1,1)$ is the
sum of its coefficients. Exactness is linear algebra.
\end{proof}

\proofstop{Write the rotated-frame bordered-minor calculation in full,
including the global sign convention. Exact tests at finitely many
dimensions confirm but do not replace the arbitrary-$(n,r)$ proof.}

\subsection{Gauge transport}

For $\widetilde F=\mu F$ with $\mu(x)=1$ and
$a=\nabla\log\mu(x)$,
\[
\widetilde g=g,\qquad
\widetilde H=H+ga^{\mathsf T}+ag^{\mathsf T},\qquad
P\widetilde HP=PHP.
\]

\begin{corollary}[gauge-channel transport]\label{cor:gauge}
The oriented shape operator and every $\sigma_r$ are fixed, while
\[
C_r(\widetilde F)-C_r(F)\in\ker\Sigma_r.
\]
For $n=3,r=2$ this is
\[
\delta\kc+\delta\kint+\delta\ks=0.
\]
\end{corollary}

\begin{remark}
This proves that reachable gauge displacements lie in the zero-sum
hyperplane. It does not prove that every vector in that hyperplane is
geometrically reachable. If $\mu(x)<0$, the chosen normal orientation
reverses and odd $\sigma_r$ change sign.
\end{remark}

\subsection{The Gaussian three-channel theorem}

For a regular implicit surface in $\R^3$,
\[
\KG=\kc+\kint+\ks.
\]
For a graph $z=h(x,y)$, with $Q_0=1+h_x^2+h_y^2$,
\[
\kc=-\frac{h_{xy}^2}{Q_0^2},\qquad
\ks=\frac{h_{xx}h_{yy}}{Q_0^2},\qquad
\kint=0.
\]
Thus the three active output-chart graph records lie on the one-dimensional
slice
\[
\kint=0,\qquad \kc+\ks=\KG
\]
inside the two-dimensional affine account plane over $\KG$.

\proofstop{The shape-operator spectrum is complete only for spectral
scalar invariants under orthogonal conjugation. Do not call
$\{\sigma_1,\ldots,\sigma_{n-1}\}$ the complete set of all
active-recharting invariants, which are governed by the larger chart-jet
quotient.}

\proofstop{Channels are invariants of a framed defining-function jet, or
canonical representatives after graph normalization. Fixing the coordinate
frame
does not fix defining-function gauge. Do not call raw channel values
gauge-invariant.}

\section{The quotient--selection separation theorem}

\begin{theorem}[two-realization quotient--selection separation]
\label{thm:separation}
The local chart and global cover realizations share the following exact
separation.
\begin{enumerate}[label=(\alph*)]
\item In the local realization, the reduced invariant $\sigma_r$ does not
determine a channel representative whenever $\ker\Sigma_r\neq0$.
Graph normalization, defining-function gauge, or a metric construction
supplies additional representative data.
\item In the global realization, the polynomial $N_k(U)$ determines an
unordered root multiset and does not by itself label one root as physical.
On the strict chamber, an oriented definite pencil selects the physical root
as its first positive boundary exit.
\item Neither selection is part of the bare invariant quotient.
\end{enumerate}
\end{theorem}

\begin{proof}
Part (a) follows from the nontrivial affine fibre of
Theorem~\ref{thm:reduction}. Polynomial coefficients are symmetric in the
roots, so they do not carry a root label. The first-exit statement is proved
in Theorem~\ref{thm:cover}. Part (c) is the conjunction.
\end{proof}

\begin{remark}
This theorem establishes a shared logical architecture. It does not provide a
functor from channel fibres to root fibres.
\end{remark}

\section{Kerr--Newman inverse-channel metric selection}

For the Kerr--Newman relation $M^2=U(S,J,Q)$, each output coordinate $i$
carries two mass--coordinate couplings and one coordinate--coordinate
coupling. Denote the coordinate--coordinate
weight by $\eta$:
\[
g_F(\eta)_{ii}
=q_i^2\left(
\Lambda_{i,\{M,j\}}^{-2}
+\Lambda_{i,\{M,k\}}^{-2}
+\eta\Lambda_{i,\{j,k\}}^{-2}
\right).
\]
Let
\[
\Wp=\{S,J,Q>0:U_S>0\}
\]
be the outer physical wedge.

\begin{theorem}[family-relative Kerr--Newman selection]
\label{thm:metricselection}
\realization.
Within the displayed weighted family:
\begin{enumerate}[label=(\roman*)]
\item each $\eta>0$ produces an interior metric and curvature pole at a
coordinate--coordinate coupling zero;
\item each $\eta<0$ loses positive definiteness near that zero;
\item $\eta=0$ is positive definite and interior-regular on $\Wp$,
while retaining the desired mass--coordinate boundary divergences.
\end{enumerate}
Thus $g_F(0)$ is the unique admissible member of this family.
\end{theorem}

\begin{proof}[Proof route]
The relevant coordinate--coordinate coupling has an interior zero. Its inverse square
enters the metric with coefficient $\eta$. Positive $\eta$ produces a
metric pole and negative $\eta$ forces the diagonal entry negative near
the zero. At $\eta=0$ that coupling is omitted. The remaining mass--coordinate
couplings are finite and nonzero in $\Wp$ and vanish only on distinguished
boundary divisors
in its closure.
\end{proof}

\proofstop{Before Theorem~\ref{thm:metricselection} bears publication
load, reproduce the exact six mass--coordinate factorizations, prove existence
of the interior coordinate--coordinate zero, and archive a symbolic certificate that the
$\eta>0$ metric pole produces a nonremovable scalar-curvature pole.}

\begin{corollary}[analytic corridor]\label{cor:interior}
\conditional.
The scalar curvature of $g_F(0)$ is analytic on $\Wp$, including the
Davies surface, and its singular support in the closure is contained in
\[
\{T=0\}\cup\{\Omega=0\}\cup\{\Phi_e=0\}.
\]
\end{corollary}

\proofstop{To replace contained in by equals, prove nonvanishing of the
leading residue on every open boundary stratum. The extremal edge depends on
the full transverse-monotonicity and Sturm certificate.}

\section{Selection realization II: algebraic normalization and reflection descent}

This section separates generic function-field assertions from pointwise real
assertions.

\subsection{Generic algebraic theorem}

Let
\[
F_0=\Q(N_1,\ldots,N_k),\quad L=F_0(m),\quad
w_i^2=m^2+N_i^2,
\]
where the $N_i$ are algebraically independent. Put
\[
E=L(w_1,\ldots,w_k),\quad s=\sum_iw_i,\quad
F=F_0(M),\quad U=m^2.
\]

\begin{theorem}[generic normalization--reflection chain]
\label{thm:genericcover}
Assume the square classes of $m^2+N_i^2$ are independent. Then:
\begin{enumerate}[label=(\roman*)]
\item $E/L$ is multiquadratic with group $H_{\mathrm{sign}}=(C_2)^k$;
\item the stabilizer of $s$ in $H_{\mathrm{sign}}$ is trivial, hence
$L(s)=E$;
\item the orbit norm
\[
\Phi_k(m,M)=
\prod_{\varepsilon\in\{\pm1\}^k}
\left(4M-\sum_i\varepsilon_iw_i\right)
\]
lies in $F_0[m,M]$, is even in $m$, and equals $N_k(U,M)$;
\item the observable-field map induces
\[
F[m]/(\Phi_k)\cong E;
\]
\item $N_k(U)$ is irreducible over $F$ and has degree
\[
\delta_k=\frac12\left(
2^k-\mathbf 1_{2\mid k}\binom{k}{k/2}
\right);
\]
\item the involution
\[
\tau(m)=-m,\qquad\tau(w_i)=w_i
\]
has fixed field $E^\tau=F(U)$, and every $w_i$ descends to $F(U)$.
\end{enumerate}
\end{theorem}

\begin{proof}[Proof route]
Square-class independence gives the multiquadratic extension. A nontrivial
sign change changes $s$, so the primitive-element stabilizer is trivial.
The orbit product is invariant under all radical signs and even under
$m\mapsto-m$. Irreducibility and the field identification follow from
the incidence map and localization over $F_0[M]$. The fixed-field statement
uses the degree-two reflection quotient.
\end{proof}

\proofstop{The companion draft says that birational affine curves have
isomorphic normalizations, which is not sufficient. If affine normalization
is claimed here, prove that the incidence curve is normal, the map to the
axial curve is finite and birational, and therefore has the normalization
universal property. The field isomorphism itself does not require that
stronger geometric sentence.}

\subsection{The pointwise real theorem}

Let $X_i,Z_i$ be Pauli matrices acting in the $i$th tensor slot and put
\[
L_0=4MI+\sum_iN_iX_i,\qquad
Z=\sum_iZ_i,\qquad
L(m)=L_0-mZ.
\]
This is the sign-conjugate form of the Nie--Parrilo--Sturmfels pencil, with
\[
\det L(m)=N_k(m^2).
\]
Define
\[
\Cham=\left\{(M,N):4M>\sum_i|N_i|\right\}.
\]

\begin{lemma}[chamber equals definiteness]\label{lem:definite}
\[
L_0\succ0\quad\Longleftrightarrow\quad(M,N)\in\Cham.
\]
\end{lemma}

\begin{proof}
The eigenvalues of $L_0$ are $4M+\sum_i\varepsilon_iN_i$; the least is
$4M-\sum_i|N_i|$.
\end{proof}

Let $X_{\mathrm{all}}=\sigma_x^{\otimes k}$. It commutes with $L_0$
and anticommutes with $Z$. In its two eigenspaces,
\[
L_0=\begin{pmatrix}A&0\\0&B\end{pmatrix},\qquad
Z=\begin{pmatrix}0&C\\C^{\mathsf T}&0\end{pmatrix},
\qquad A,B\succ0.
\]
Set $T=A^{-1/2}CB^{-1/2}$.

\begin{theorem}[spectral chamber and physical selection]
\label{thm:cover}
\established.
Let
\[
D_k=
\begin{cases}
2^k,&k\text{ odd},\\
2^k-\binom{k}{k/2},&k\text{ even},
\end{cases}
\qquad\delta_k=D_k/2.
\]
For every real $(M,N)\in\Cham$, including zero or repeated charges:
\begin{enumerate}[label=(\roman*)]
\item $\rank C=\delta_k$ and
\[
N_k(U)=\det(L_0)
\prod_{j=1}^{\delta_k}
\left(1-U\sigma_j(T)^2\right);
\]
\item every root of $N_k$ is real and strictly positive;
\item each sign vector with positive imbalance
$a_\varepsilon=\sum_i\varepsilon_i>0$ crosses the level $4M$
exactly once and transversely, while zero- and negative-imbalance sheets do
not cross;
\item the all-plus root $U_{\mathrm{phys}}$ is the unique least root;
\item
\[
U_{\mathrm{phys}}
=\lambda_{\max}\left(
L_0^{-1/2}ZL_0^{-1/2}
\right)^{-2}
=\left[
\max_{v\neq0}\frac{v^{\mathsf T}Zv}
{v^{\mathsf T}L_0v}
\right]^{-2};
\]
\item if $N_k(U)=\sum_{r=0}^{\delta_k}c_rU^r$, then
$(-1)^rc_r>0$ for every $r$.
\end{enumerate}
\end{theorem}

\begin{proof}
Parity makes the normalized symmetric pencil direction off-diagonal. Its
nonzero eigenvalues are the pairs $\pm\sigma_j(T)$, yielding the
factorization and positive roots.

Every signed sheet begins below $4M$ at $U=0$. A positive-imbalance sheet
tends to $+\infty$, so crosses at least once. The number of such sheets is
$\delta_k=\deg N_k$; their crossings already saturate total degree
counted with multiplicity. Hence each is unique and transverse and there
are no further crossings.

At a shadow crossing,
\[
\sum_iw_i(U)
=4M+2\sum_{\varepsilon_i=-1}w_i(U)
>4M=\sum_iw_i(U_{\mathrm{phys}}).
\]
Strict monotonicity of the all-plus sum gives $U>U_{\mathrm{phys}}$.
The generalized-eigenvalue formula follows, and coefficient alternation is
the expansion of the positive singular-value product.
\end{proof}

\begin{corollary}[four-channel global simplicity]
For $k=4$, if $N_i^2\neq N_j^2$ for $i\neq j$, then all five roots are
simple at every strict-chamber point. The physical root is first, and larger
$|N_i|$ gives the later associated single-minus shadow root.
\end{corollary}

\begin{proof}
The positive-imbalance sheets are the all-plus sheet and four single-minus
sheets. Each crossing is transverse. Two single-minus sheets collide exactly
when $w_i=w_j$, equivalently $N_i^2=N_j^2$.
\end{proof}

\begin{corollary}[physical analytic section]
The map $(M,N)\mapsto U_{\mathrm{phys}}(M,N)$ is real-analytic on the
strict chamber and is fixed by continuation around every loop contained in
that chamber.
\end{corollary}

\begin{proof}
The physical factor has nonzero $U$-derivative everywhere in the chamber, so
the implicit-function theorem gives local analyticity. Global uniqueness
glues the local sections and prevents nontrivial continuation.
\end{proof}

\proofstop{State explicitly that this is the sign-conjugate NPS pencil.
Do not call the sign convention character-for-character identical to the
source.}

\proofstop{The critical-sheet discriminant mechanism is absent for
strict-chamber axial crossings, but it remains a genuine mechanism off
chamber and in the complex/global discriminant. Do not erase it globally.}

\section{Four-charge arithmetic consequences}

Let $F=\Q(M,N_1,N_2,N_3,N_4)$. The generic $N_4(U)$ has degree five.

\begin{theorem}[$S_5$ obstruction]\label{thm:S5}
\conditional.
Assuming the exact Dedekind certificates described in the companion
calculation, $N_4$ is irreducible over $F$ with generic Galois group $S_5$.
Consequently the generic seed mass-square is not expressible by radicals
over the observable field.
\end{theorem}

\begin{proof}[Certificate route]
At one rational specialization, exact irreducibility makes the specialized
group transitive, while a good-prime factorization of cycle type $(2,3)$
forces that transitive quintic group to be $S_5$. The separable specialized
group embeds into the generic group, which is a subgroup of $S_5$.
\end{proof}

\proofstop{Include the specialized quintic, irreducibility certificate,
good-prime verification, and modular factorization in the paper or a
permanent machine-readable supplement. A prose report that a run occurred
is not an auditable certificate.}

\begin{corollary}[totally real arithmetic, correctly scoped]
Every strict-chamber specialization has five positive real roots counted
with multiplicity. Any rational or algebraic chamber specialization that is
separately certified irreducible with Galois group $S_5$ defines a totally
real $S_5$ quintic field.
\end{corollary}

\begin{remark}[specialization guardrail]
Generic $S_5$ does not persist at every chamber point. At
$(M,N_1,N_2,N_3,N_4)=(2,1,1,1,1)$,
\[
N_4(U)=2^{30}(3-U)(15-U)^4.
\]
The physical root is simple while four shadow sheets collide.
\end{remark}

\begin{proposition}[ordered horizon reconstruction]
\conditional.
Under the generic descent and square-class hypotheses, signed radicals lie
in $F(U)$, the oriented axis field is obtained by adjoining $m$, and the
ordered horizon field is generated by either normalized horizon entropy.
The rational entropy product reconstructs its conjugate.
\end{proposition}

\proofstop{Keep unconditional field identification separate from the
proposed degree-$20$ count. The latter still requires an independent
square-class gate and a good-specialization argument.}

\section{Involution, parity, and local curvature valuations}

The cover theorem uses an involution commuting with the base matrix and
anticommuting with the pencil direction. The inverse-channel boundary theorem uses a
coordinate reflection fixing a metric germ. In both cases odd data vanishes,
but the actions and outputs remain distinct.

\begin{lemma}[fixed-germ parity]
If a scalar invariant $R(x,y)$ is constructed naturally from a metric germ
whose components are even in $x$, then $R(-x,y)=R(x,y)$. Its Laurent
expansion contains no odd powers of $x$.
\end{lemma}

\begin{proof}
The reflection is an isometry of the germ. Naturality of scalar curvature
gives the result.
\end{proof}

\begin{proposition}[generic quadratic collapse]\label{prop:generic}
Let a diagonal metric satisfy
\[
g_{xx}=Ax^2+O(x^3),\qquad
g_{\alpha\alpha}
=P_{\alpha0}+P_{\alpha1}x+O(x^2),
\quad AP_{\alpha0}\neq0.
\]
Then
\[
R=
\frac1A\sum_{\alpha=1}^{m}
\frac{P_{\alpha1}}{P_{\alpha0}}x^{-3}
+O(x^{-2}).
\]
The pole has exact order three only if the displayed coefficient is nonzero.
\end{proposition}

\begin{theorem}[pure parity-fixed inverse-channel normal form]
\label{thm:paritynormal}
For $m\geq1$, $B>0$, and $A_\alpha>0$,
\[
ds^2=Bx^2dx^2+x^{-2}\sum_{\alpha=1}^{m}
A_\alpha dy_\alpha^2
\]
has scalar curvature
\[
R=-\frac{m(m+5)}B x^{-4}.
\]
The coefficient is independent of all transverse amplitudes.
\end{theorem}

\begin{proof}
Set $\rho=\frac12\sqrt B\,x^2$. The metric becomes a warped product
$d\rho^2+w(\rho)^2h_{\mathrm{flat}}$ with
$w\propto\rho^{-1/2}$. The flat-fibre warped-product formula gives
\[
R=-2m\frac{w''}{w}
-m(m-1)\left(\frac{w'}w\right)^2
=-\frac{m(m+5)}{4\rho^2}
=-\frac{m(m+5)}{Bx^4}.
\]
\end{proof}

\begin{corollary}[asymptotic fixed face]
Assume a diagonal, even, $C^2$ polyhomogeneous germ for which
\[
g_{xx}=B(y)x^2+O(x^4),\qquad
g_{\alpha\alpha}=A_\alpha(y)x^{-2}+O(1),
\]
with positive $B,A_\alpha$ and uniform control of the stated remainders
and transverse derivatives. Then
\[
R=-\frac{m(m+5)}{B(y)}x^{-4}+O(x^{-2}).
\]
\end{corollary}

\proofstop{Write the full asymptotic curvature expansion proving that
transverse derivatives do not create another $x^{-4}$ contribution. If the
needed uniform/product hypotheses are stronger than stated, add them.}

\subsection{Corners}

The current general codimension-two statement is retained only in a
monomial/polyhomogeneous form.

\begin{proposition}[conditional corner vertex rule]\label{prop:corner}
\conditional.
Suppose a diagonal metric has specified monomial weights
\[
g_i=h_i(s)x^{p_i}y^{q_i}u_i(x,y,s),
\]
where each $u_i$ is a nonvanishing polyhomogeneous unit and the polar
coefficients below are nonzero. Then the polar part of scalar curvature is
supported on the candidate vertices
\[
V_1=(-(p_1+2),-q_1),\quad
V_2=(-p_2,-(q_2+2)),\quad
V_0=(-p_0,-q_0).
\]
Along $x=\rho\epsilon^{a_x}$, $y=\epsilon^{a_y}$, the pole order is
the support function of the surviving vertices.
\end{proposition}

\proofstop{Single-face pole orders and residue weights alone do not
determine all mixed corner monomials. Include the closed symbolic identity
for the coefficients of $V_0,V_1,V_2$, derive it from the diagonal
scalar-curvature formula, and state every cancellation condition. Do not
extrapolate to codimension at least three.}

\proofstop{Deck parity $m\mapsto-m$, spin reflection $J\mapsto-J$,
charge reflection $Q\mapsto-Q$, and defining-function gauge are distinct
actions. Their common $C_2$ representation pattern is a bridge, not an
identification.}

\section{Kerr--Newman complementarity as boundary realization}

\begin{theorem}[open-stratum inverse-channel complementarity]
\label{thm:344}
\conditional.
For the selected metric $g_F(0)$:
\begin{enumerate}[label=(\roman*)]
\item scalar curvature is analytic in the open outer wedge, including on
the Davies surface;
\item on the open generic extremal stratum $T=0$, away from its
intersections with reflection faces, curvature has exact pole order three;
\item on the open reflection-fixed strata $\Omega=0$ and $\Phi_e=0$,
away from further intersections, curvature has exact pole order four and
leading coefficient $-14/B$;
\item the verified Schwarzschild intersection
$\Omega=\Phi_e=0$ has a direction-dependent Newton wedge.
\end{enumerate}
\end{theorem}

\begin{proof}[Dependency proof]
Interior regularity is Corollary~\ref{cor:interior}. The extremal germ is
of the form in Proposition~\ref{prop:generic}; its separately certified
nonzero coefficient gives exact order three. Each open reflection face is
the $m=2$ case of Theorem~\ref{thm:paritynormal}. The Schwarzschild
calculation is the specialized instance of Proposition~\ref{prop:corner}.
\end{proof}

\proofstop{Re-run and archive the factorization, interior-regularity,
extremal noncancellation/Sturm, reflection-residue, and corner certificates.
The intersections $T\cap\Omega$ and $T\cap\Phi_e$ are not
classified and must be excluded or separately proved.}

\proofstop{For a path $Q=\epsilon$, $J=\rho\epsilon^a$ to approach
the Schwarzschild corner one needs $a>0$. The endpoint $a=0$ approaches only
the charge face with nonzero $J$ and must not be described as a corner path.}

\begin{remark}[scope]
The $3/4/4$ law is portable only conditionally: another diagonal
inverse-channel surface must first be proved to have the corresponding local
germs and nonzero coefficients. It is not a theorem about every framed
constraint hypersurface.
\end{remark}

\section{Global continuation and arithmetic periods remain companions}

\subsection{Self-glue monodromy}

For a self-glued algebraic relation, solving different output coordinates produces
different branched covers. There is generally no single Galois or monodromy
group attached to the unlabelled surface independently of the extraction
coordinate.

\begin{proposition}[local and global symmetry are independent]
\bridge.
After lifting local chart-labelled data to every analytically continued
solution sheet, global monodromy acts on the sheet index while local chart
permutation acts on the coordinate index. When analytic continuation preserves
the coordinate labels, the two actions commute; they are not canonically
identified.
\end{proposition}

\proofstop{The companion self-glue draft needs substantial repair before
its general wreath laws are imported. In particular, $C_m\wr C_m$ is
usually nonabelian and is not cyclic; the general gcd-layer kernel needs a
Kummer or branch-cycle proof; a full directed-wreath realization requires independent
within-block discriminant loci, not merely distinct block labels; and all
quantifiers must respect the stated exceptional divisors.}

\subsection{Corrected curvature constants}

The elliptic curvature constants concern global integration cycles,
third-kind differentials, an isogeny, and principal-value residue jumps.
They are not used to prove the local orbit, metric-selection, or axial-cover
theorems.

\begin{theorem}[primary surface-to-link close-off]\label{thm:surface-link}
\established.
For
\[
S_s=\{D^2+sDS+P^2=3\},\qquad s\ne0,
\]
let $\Gamma_s^+$ be either component of the spherical link of the
asymptotic cone. Then
\[
\int_{S_s}K_G\,dA=-2L(\Gamma_s^+).
\]
At $s=1$, exact Legendre reduction of $L(\Gamma_1^+)$ gives
\[
\int_{S_1}K_G\,dA
=-2^{7/4}\left[
(3+2\sqrt2)\Pi\!\left(\frac{4-3\sqrt2}{8};
\frac{2-\sqrt2}{4}\right)
-(2+2\sqrt2)K\!\left(\frac{2-\sqrt2}{4}\right)
\right]
=I_{\rm primary}.
\]
\end{theorem}

\begin{proof}
Orthogonally diagonalize the surface to
$ax^2-by^2+z^2=3$, where
$a=(1+\sqrt{1+s^2})/2$ and
$b=(\sqrt{1+s^2}-1)/2$.  On this cylinder the one-form
\[
\eta=-\frac{2ab\,y}{\sqrt q\,(a^2x^2+z^2)}
(x\,dz-z\,dx),\qquad q=\lVert\nabla F\rVert^2,
\]
is globally regular and satisfies $d\eta=K_GdA$.  In the global
hyperbolic parametrization its coefficient tends at the two oriented
ends to
\[
-\frac{\sqrt{ab}}{p(t)\sqrt{p(t)+b}},
\qquad p(t)=a\cos^2t+\sin^2t.
\]
The reciprocal tangent substitution
\[
\tan t=\sqrt{\frac{a+b}{1+b}}\cot\theta
\]
identifies the integral of this density with the length of one component
of the spherical asymptotic link.  Stokes' theorem therefore gives one
copy from each end and proves the displayed factor $-2$.  The final
formula is the established link-to-Legendre reduction.
\end{proof}

\proofstop{Theorem~\ref{thm:surface-link} closes the primary scalar
surface-to-period map. It does not identify the Galois-conjugate
principal-value constant with a real region or relative cycle of the
original DBP surface. A chain-level pushforward of the original curvature
two-form to the elliptic differential would refine the primary statement,
but is no longer needed to establish its scalar value; the dual surface
cycle remains genuinely open.}

\begin{remark}
The primary and dual channels now have different geometric status.  The
primary period is the proved total curvature of $S_1$.  For the explicitly
defined dual third-kind differential, the weighted residue is four and the
two oriented indentations differ from the principal value by $\pm4\pi i$,
hence from one another by $8\pi i$; its interpretation in the original DBP
realization still requires a new differential-and-cycle theorem.
\end{remark}

\section{What the paper absorbs and what it retires}

The following results are internal to the main spine:
\begin{enumerate}
\item passive versus active output-chart correction;
\item the Active Recharting Orbit Theorem;
\item bordered-minor channel reduction and gauge transport;
\item the three-channel Gaussian specialization;
\item the family-relative Kerr--Newman metric selection;
\item generic and parity-fixed local curvature normal forms;
\item the generic normalization--reflection chain;
\item the axial definite-pencil, sheet-saturation, and least-root
selection theorem.
\end{enumerate}

The former curvature-orbit correction memo and canonical reduction note can
be retired into the present text. The implementation-focused three-channel
specification becomes a verification appendix. Detailed self-glue
monodromy, elliptic periods, rotating horizon reconstruction, and full
arithmetic certificates remain companion work.

\section{The remaining unification theorem}

\begin{conjecture}[intrinsic orbit--quotient--selection principle]
There is a class of systems equipped with
\[
(\text{carrier},\text{ symmetry},\text{ quotient},
\text{ real chamber},\text{ selected section},\text{ wall data})
\]
and morphisms preserving these data, such that:
\begin{enumerate}[label=(\roman*)]
\item the local active-recharting systems and global horizon-cover systems are
objects of the class;
\item invariant observables factor through the quotient;
\item physical observables require the selected section;
\item selected sections avoid the interior branch locus;
\item boundary degenerations are typed functorially by stabiliser and wall
data.
\end{enumerate}
\end{conjecture}

\proofstop{This is not a theorem until the data are minimized, native
morphisms and equivalence are defined, closure is proved, and a third
independent realization is exhibited. It must not be used to prove any
preceding result.}

\section{Conclusion}

The unified framework is a local-to-global theorem chain. The active
chart-jet orbit is the local carrier; elementary shape curvature is a reduced
projection; channel data is a framed account. A physical realization then
requires selection: in Kerr--Newman geometry, positivity and interior
regularity select a metric-family member; in the axial cover, definiteness
and orientation select the first boundary root. Involutions remove odd data
and organize fixed strata, while local metric germs determine boundary
valuations.

What has been unified is the logical architecture and dependency of proved
results. What has not been proved is a categorical equivalence of the
realizations. Keeping that boundary explicit is what makes the synthesis
auditable.

\appendix

\section{Source absorption map}

\begin{center}
\small
\begin{tabular}{>{\raggedright\arraybackslash}p{0.30\textwidth}
>{\raggedright\arraybackslash}p{0.42\textwidth}
>{\raggedright\arraybackslash}p{0.18\textwidth}}
\toprule
Source manuscript & Disposition & Status\\
\midrule
Active role-jet orbit calculus &
Foundational local theorem, with faithfulness correction & Absorbed\\
Canonical Invariant Reduction and Gauge Transport &
Channel polynomial and gauge section & Absorbed\\
Theorem 8.1 curvature correction &
Passive/active distinction and value-plus-jet guardrail & Retired memo\\
Three-channel $K_G$ specifications &
Gaussian specialization and verification appendix & Absorbed in part\\
Kerr--Newman inverse-channel metric &
Selection and complementarity realization & Conditional\\
Parity-fixed normal forms &
Local valuation section & Absorbed; corner proof owed\\
Galois horizon cover and R5--R7 entries &
Global cover module & Conditional arithmetic\\
Self-glue monodromy & Global continuation companion & Not load-bearing\\
Corrected curvature constants &
Primary surface-to-link theorem absorbed; dual arithmetic-period companion &
Primary established; dual geometric cycle open\\
\bottomrule
\end{tabular}
\end{center}

\section{Notation collision ledger}

\begin{center}
\begin{tabular}{lll}
\toprule
Concept & Unified notation & Avoid\\
\midrule
mass square & $U=m^2$ & metric-family $u$\\
coordinate--coordinate metric weight & $\eta$ & mass-square $u$\\
channel polynomial variable & $v$ & mass-square $U$\\
gradient norm & $q=g^{\mathsf T}g$ & charge $Q$ without context\\
active chart group & $G_{\mathrm{act}}=S_3$ & global monodromy group\\
sign group & $H_{\mathrm{sign}}=(C_2)^k$ & axial deck group\\
axial deck group & $G_{\mathrm{deck}}=\langle\tau\rangle$ & charge reflection\\
mass chamber & $\Cham$ & outer KN wedge $\Wp$\\
channel tuple & $(\kc,\kint,\ks)$ & alternating tuple orders\\
\bottomrule
\end{tabular}
\end{center}

\section{Bibliographic placeholders}

\begin{thebibliography}{99}
\bibitem{NPS}
J. Nie, P. A. Parrilo, and B. Sturmfels,
\emph{Semidefinite Representation of the $k$-Ellipse},
IMA Volumes in Mathematics and its Applications 146 (2008), 117--132.

\bibitem{PV}
D. Plaumann and C. Vinzant,
\emph{Determinantal Representations of Hyperbolic Plane Curves:
An Elementary Approach},
J. Symbolic Computation 57 (2013), 98--107.

\bibitem{Orbit}
W. Lloyd,
\emph{The Active Role-Jet Orbit Calculus for Constraint-Surface Roles},
companion manuscript.

\bibitem{KN}
W. Lloyd,
\emph{Exact Role-Divisor Retrodiction in Kerr--Newman Thermodynamics
via an Inverse-Channel Complementarity Metric},
companion manuscript.

\bibitem{PFC}
W. Lloyd,
\emph{Parity-Fixed Inverse-Channel Normal Forms and Universal Curvature
Blow-Up},
companion manuscript.

\bibitem{Galois}
W. Lloyd,
\emph{Galois Theory of the Horizon Cover},
companion manuscript.
\end{thebibliography}

\end{document}
